\documentclass[12pt]{article}
\usepackage[utf8]{inputenc}
\usepackage[T1]{fontenc}
\usepackage{hyperref}
\title{Tetela Language Resources for K-12 Curriculum Development}
\author{Compiled by Claude Sonnet 5 (claude-sonnet-5) for Andr\'as Kornai}
\date{}
\begin{document}
\maketitle

\section{Introduction}
Tetela (Otetela, Kitetela) is a Bantu language of northern Kasai-Oriental, DR Congo. Speaker estimates range from about 760,000 (Wikipedia) to about 1 million (Omniglot). This search found a historical dictionary, a catalogued grammar and a reported Bible translation. It found no school materials.

\section{Primary Educational Resources}
None found in this search.

\section{Linguistic and Reference Materials}
\begin{itemize}
\item Otetela and English dictionary (Episcopal Congo Mission, Wembo-Nyama, 1925), digitized at SOAS: \url{https://digital.soas.ac.uk/AA00000710/00001/citation}
\item Glottolog reference entry for Omatete, Alonge-Dikonda 1985 (grammar). The type of work is uncertain, and I did not check the work itself: \url{https://glottolog.org/resource/reference/id/86282}
\item Wikipedia article on the Tetela language, including its noun-class system: \url{https://en.wikipedia.org/wiki/Tetela_language}
\item Wikipedia article on the Tetela people: \url{https://en.wikipedia.org/wiki/Tetela_people}
\end{itemize}

\section{Religious and Community Resources}
\begin{itemize}
\item Omniglot page on Tetela. It states that a Bible translation exists and gives a sample passage, but I did not verify availability: \url{https://www.omniglot.com/writing/tetela.htm}
\end{itemize}

\section{Digital and Audio Resources}
None found.

\section{Curriculum Development Context}
Tetela has a large speaker base but few documented educational resources. Existing materials are mostly historical or religious. A curriculum would need new primers and readers. The 1925 dictionary and the Bible text could serve as source material, after checking their orthography.

\end{document}